% 【本文件写什么】独立子 crate，无 api/impl 目录树：driver-impl
% - 职责：板级驱动 impl：两个 QEMU profile 与两个真机 profile。
% - 源码：os/components/wateros-driver/driver-impl/src/
% - 对外暴露的函数/类型、在聚合 lib.rs 中的重导出方式
% - 实现要点、feature，若有、与其它子 crate 的边界
% 【事实来源】
% - os/components/wateros-driver/driver-impl/、父组件 src/lib.rs

\subsection{driver-impl}

\code{driver-impl}承载平台相关实现：QEMU RISC-V、QEMU LoongArch64 virt、VisionFive 2/JH7110 与 Loongson 2K1000LA。聚合\code{lib.rs}通过互斥 feature 返回对应\code{MachineDriver}；未选择平台会显式失败，不提供 dummy 回退。各 RISC-V 与 LoongArch 平台按需导出\code{uart}子模块供早期串口与\code{klog}使用。

平台impl职责边界：解析DTB或PCI配置空间、实例化各子系统\code{impl-virtio-*}驱动、调用\code{register\_*}、同步\code{devfs}，以及推导物理RAM上界供MM恒等映射。不包含文件系统挂载或socket系统调用逻辑。重复\code{init\_after\_boot}幂等：第二次调用记录\code{WARN}并忽略，避免重复注册。

% 【本文件写什么】impl 实现层：qemu-loongarch64-virt
% - 定位：QEMU LoongArch64 virt 板级组合。
% - crate 路径：os/components/wateros-driver/driver-impl/impl-qemu-loongarch64-virt/src/
% - 如何实现对应 api-v0 trait：关键类型、状态机、数据结构
% - 算法或硬件交互流程，可用伪代码/流程图
% - 本 impl 启用的 Cargo [features] 与 arch feature，impl-riscv64 等
% - 与其它 impl 变体的差异、切换 feature 时的影响
% - 性能/内存假设与已知限制
% 【事实来源】
% - 上述 crate 源码与 Cargo.toml
% - os/feature-tree.txt 中指向本 impl 的 feature 链

\subsubsection{impl-qemu-loongarch64-virt}

\paragraph{与RISC-V路径差异。}块/网卡不走DTB virtio-mmio扫描，而在\code{init\_after\_boot\_inner}中经PCIe ECAM枚举\code{virtio-blk}/\code{virtio-net}。\code{pci::find\_config\_base}从DTB PCI节点或回退常量解析配置空间基址；\code{probe\_virtio\_blk\_pci}/\code{probe\_virtio\_net\_pci}分别调用块/网\code{impl-virtio-pci}的\code{probe\_first\_from\_ecam}。

\paragraph{引导与RAM。}\code{init\_when\_boot}保存DTB并\code{uart::init\_early\_default\_uart}初始化\code{0x1FE0\_01E0} NS16550。\code{physical\_ram\_end\_exclusive}优先DTB\code{memory@*}，失败回退\code{0xb000\_0000}以匹配QEMU\code{virt -m 1G}高段布局；内核自 \code{0x9000\_0000} 启动，避免把MMIO空洞当作RAM。

\paragraph{注册与devfs。}块设备成功时可选\code{BlockCacheManager::wrap}并\code{register\_block\_device}，记录\code{VirtioPciProbeInfo}到\code{VIRTIO\_BLK\_PCI}。网卡\code{register\_network\_device}并记录MAC与BDF。字符路径调用\code{uart::register\_default\_uart}，全局 \code{QemuLoongArch64Uart16550}，与 \code{register\_builtin\_character\_devices}。末尾\code{devfs\_impl::refresh}暴露\code{/dev/vblk*}等节点。

\begin{rustcode}
pub fn init_when_boot(dtb_pa: usize) {
    DTB_BASE_ADDR.store(dtb_pa, Ordering::Release);
    uart::init_early_default_uart();
}

pub fn physical_ram_end_exclusive() -> usize {
    // DTB memory@* 或回退 0xb000_0000
    0xb000_0000
}
\end{rustcode}

\paragraph{导出接口。}除与RISC-V相同的\code{init\_after\_boot}/\code{physical\_ram\_end\_exclusive}/\code{test}外，\code{uart::with\_default\_uart}供\code{klog}与早期打印在已初始化UART上执行闭包。重复\code{init\_after\_boot}同样幂等忽略。

\paragraph{限制。}仅扫描PCI bus0首个块/网设备；DTB扫描不用于virtio绑定，\code{scan\_device\_info} 非本路径主流程。中断未启用，IO与收包依赖轮询。多实例devfs与syscall策略仍以首设备为主。

% 【本文件写什么】impl 实现层：qemu-riscv64-opensbi
%   - 定位：QEMU RISC-V + OpenSBI 板级组合。
%   - crate 路径：os/components/wateros-driver/driver-impl/impl-qemu-riscv64-opensbi/src/
%   - 如何实现对应 api-v0 trait：关键类型、状态机、数据结构
%   - 算法或硬件交互流程（可用伪代码/流程图）
%   - 本 impl 启用的 Cargo [features] 与 arch feature（impl-riscv64 等）
%   - 与其它 impl 变体的差异、切换 feature 时的影响
%   - 性能/内存假设与已知限制
% 【事实来源】
%   - 上述 crate 源码与 Cargo.toml
%   - os/feature-tree.txt 中指向本 impl 的 feature 链

\subsubsection{impl-qemu-riscv64-opensbi}

\paragraph{引导状态。}\code{init\_when\_boot}将DTB物理地址存入\code{DTB\_BASE\_ADDR}。\code{physical\_ram\_end\_exclusive}遍历DTB\code{memory@*}节点\code{reg}，取\code{base>=0x8000\_0000}的最大末端；失败回退\code{wateros-base-config::QEMU\_VIRT\_PHYS\_RAM\_END}。

\paragraph{DTB扫描。}\code{scan\_device\_info}清空并重建\code{DEVICE\_INFOS}：解析\code{compatible}列表、首段MMIO、IRQ与virtio魔数/\code{device\_id}探测类型。对\code{virtio,mmio}节点打诊断日志（魔数、\code{device\_id}、区域大小）。

\paragraph{设备注册流程。}\code{init\_after\_boot\_inner}顺序为：打印各子系统\code{supported\_devices}目录、\code{scan\_device\_info}、\code{probe\_character\_devices}、\code{probe\_virtio\_blk\_and\_collect\_unsupported}。字符路径：DTB NS16550节点注册\code{SerialPortCharacterDevice}；若无节点则回退\code{QemuVirtUart16550::qemu\_virt\_default()}（\code{0x1000\_0000}）；最后\code{register\_builtin\_character\_devices}。块/网路径：对认领的virtio-mmio节点调用\code{VirtioBlkDevice::from\_mmio}/\code{VirtioNetDevice::from\_mmio}，可选\code{BlockCacheManager::wrap}后\code{register\_*}。

\paragraph{devfs与自检。}\code{sync\_devfs}将未绑定virtio节点路径交给\code{devfs::set\_dt\_unsupported\_paths}并\code{refresh}。\code{virtio\_blk\_probe\_test}对首块读LBA0验证IO。\code{test()}在\code{init\_after\_boot}完成后转储\code{DeviceInfo}与devfs节点列表。重复\code{init\_after\_boot}由\code{INIT\_AFTER\_BOOT\_DONE}原子标志忽略。

\begin{rustcode}
fn probe_virtio_device_type(mmio: MmioRegion) -> DeviceType {
    let magic = mmio_read32(mmio.base, 0);
    let device_id = mmio_read32(mmio.base, 2);
    if magic != 0x74726976 {
        return DeviceType::Unknown;
    }
    match device_id {
        2 => DeviceType::Block,
        1 => DeviceType::Network,
        _ => DeviceType::Unknown,
    }
}
\end{rustcode}

\paragraph{限制。}\code{DeviceInfo.irq}已解析但未挂接PLIC处理；多盘/多网卡仅全部注册，根FS与\code{stack::init}仍偏向首设备。块设备写路径依赖virtio-blk实现，缓存为写穿。

